\documentclass{article}
\usepackage{graphicx} % Required for inserting images

\usepackage{revy}
\usepackage[utf8]{inputenc}
\usepackage[T1]{fontenc}
\usepackage[danish]{babel}
\usepackage{amsmath}
\usepackage{amssymb}   % Required for blackboard bold (\mathbb) and fraktur (\mathfrak)
\usepackage{mathrsfs}

\revyname{MatematikRevyen}
\revyyear{2026}
\version{1}
\eta{$3$ minutter}
\status{Faerdig}

\title{InDRUKtionsbeviset}
\author{Daniel '23}

\begin{document}

\maketitle

\begin{roles}
    \role{EH}[Ernst Hansen]
    \role{S1}[Studerende]
    \role{S2}[Studerende]
\end{roles}

\begin{sketch}
    \scene{Scenen er sat op som en forelæsning. Til venstre er der en tavle, og til højre nogle studerende, der sidder og venter på Ernst. Skechten begynder med "The Imperial March" (fra Star Wars), som spilles. Ernst kommer ind på scenen og går hentil tavlen}

    \says{EH} JA, VI GÅR IGANG IGEN!
    \scene{The Imperial March dør}
    \says{EH} Yeah, i sidste uge pillede vi tøjet af og indgik forespillet til det evntyrlige begreb, som netop er inDRUKtion, men vi fik ALDRIG serveret bevistet for dette. Så derfor indrager vi samtykligt samleje med beviset til indruktion, right NOW! Beviset for indruktion siger simplethen, at man ALTID kan drikke én øl mere. Læg mærke til, at vi her nøje og anstændigt har valgt vores ALCOHOL Of Choice til at være øl. Denne lille enorme detalje gør beviset lettilgængligt for selv den mest uskyldige rus. Sure, man kan også gennemføre beviset med et andet Alcohol Of Choice, men det bevis overlades til publikum. Lad os påbegynde beviset! Altså hvad er indruktionsstarten. WELLLL, det hænder naturligvis således, at man altid kan drikke blot én øl.

    \scene{Ernst tegner et $1\cdot $øl på tavlen (Han skriver på tavlen med sin voldsome Ernst type beat måde}

    \says{EH} Det er en eksistenssætning, som vi beviste før pausen. Næste stop er indruktionsskridtet. Nu antager vi, at man kan drikke $n$ øl. 
    
    \scene{Ernst tegner $n\cdot$øl på tavlen under det ovenstående}

     \says{EH} Husk eksistenssætningen påstår, at man altid kan drikke én øl. MEN vi kunne i forvejen per antagelse og obvious commom sense drikke $n$ øl. Og nu yderligere én til. Jamen så kan vi altså drikke $n+1$ øl. WELLLL, så er verdens vigtigste sætning hermed bevist. Man kan ALTID drikke én øl mere. Nogle spørgsmål?

    \says{S1} \act{med poten i vejret} Men Ernst H, sammenbrydes beviset, hvis man overvejer, om det gælder for overtælleligt mange øl?"

    \says{EH} WELLLL, hvis man skal drikke overtælligt mange øl, så skal man FANDME drikke hurtigt! Men ikke desto mindre er det muligt. Men så bevæger vi os over i den dybe ende af bassinet. Altså, hvad der kaldes "kontinuert druk".

    \says{S1} Lige et tillæsgsspørgsmål hvornår skal vi have om det, Mr. Ernst?

    \says{EH} WEELLL, I skal have om det i næste blok, hvor I skal have Likørintergral og Ølteori, som min søn ulDRIK er forelæser i. Andre spørgsmål?
    
    \says{S2} \act{Også med poten i verjet} Altså man kan altid drikke én øl mere. Jeg troede, det hed noget med 3 GT'er per akt.
    
    \says{EH} Nååårh ja, der er jo de tre GT'er: Guld Tuborg, Gin og Tonic og Galoisteori. Man så skal I tage Alkobra 3 med Ian Kimming. Anyways, i næste uge vil vi bevise sætningen: Hvis man afbilder øl over i den tomme mængde, medfører det, at man selv bliver afbildet over i den fulde mængde.
\end{sketch} 


\end{document}
